\scanpage{133}
\mspage{66}

telle que $D_J\ne\varnothing$, \inserted{(i.e. $\unsure{\bar D_J|D_J}=\bigcap_{j\in J}D_j$ \unsure{soit} $\ne\varnothing$)} on trouve sur $I$ une structure de \deleted{\unsure{solide}} \inserted{maquette} simpliciale --- quelle est sa signification topologique? Les \unsure{composantes|comp.} connexes correspondent aux comp. connexes de $D=X\setminus U$ --- ou encore : celles de $\deleted{\unsure{U_\infty}}\hat X\setminus D$, où $\hat X$ est le \inserted{\unsure{l'ombre}} complété de $X$ le long de $D$ --- et $\hat X\setminus D$ a un caractère invariant p.r. à la ``compactification'' choisie $X$ de $U$. S'il était vrai que les $D_J$ \inserted{non vides} sont $\unsure{\infty}$-\unsure{connexes} (ce qui bien sûr n'est \uline{\unsure{pas}} le cas, sauf pour $D_J$ réduit : un point!) on pourrait penser que la topologie de notre \deleted{\unsure{solide}} \inserted{maquette} simpliciale exprime le type d'homotopie de $U_\infty$.

Considérons maintenant le ``localisé'' \deleted{\unsure{de}} (au sens des \unsure{ombres}, ou en un sens étale convenable) de $X$ le long de $D_J$, soit $\deleted{\unsure{\hat X_J}}X_J$, et soit $U_J$ \inserted{$=X_J\setminus D_J$,} la trace de $\unsure{U}$ sur $X_J$. Considérons $\mathcal E_{U_J}$ (de façon précise, $\mathcal E_U^{\bar L}$, si on localise \deleted{p.r. à} étalement), on aura une suite exacte
\[
\tag{37}
\unsure{T^J|\bar T^J}\unsure{\hookrightarrow|\to}\mathcal E_{U_J}^{\bar L}\longrightarrow\mathcal E_{D_J}^{\deleted{\unsure{\tilde D_J}}}\longrightarrow1
\]
--- de façon précise, choisissant un rev. universel $\tilde D_J$ de $D_J$, on trouvera un rev. universel \inserted{(par image inverse)} $\tilde X_J$ de $X_J$, d'où d'ailleurs $\tilde U_J$ et $\tilde U_J\longrightarrow U$ ; \deleted{\unsure{et on}} on désigne par $\tilde L_J$ le corps des fonctions de $\tilde X_J\unsure{\mathrel{\text{---}}|=|\ldots}\tilde U_J$, c'est une ext. \inserted{\unsure{alg.}} de $L$, et

\margin{il vaut mieux travailler avec les ``\unsure{tubulures} étales''\ldots}

\margin{il faut choisir $\tilde D_J$ ``au dessus de $\bar K$''\ldots}

\margin{plutôt $T^J$}

\margin{\uline{NB} Il faut partir de $\unsure{U_{J\bar K}}$, \unsure{si on veut}, d'un rev. universel $\unsure{\tilde V_0}$ de $\unsure{U_{J\alpha}}$ (qu'on va \unsure{multiplier} au dessus de $\bar K$ bien sûr), d'où $\unsure{\tilde X_{J\alpha}}$, $\unsure{\tilde D_{J\alpha}}$.}

\margin{132}
